% Source: physical PDF pages 487--491; printed pages 464--468.
% Section 23.8 begins at the heading on physical p. 487 and ends before
% Appendix A on physical p. 491.
% Figures: none. Source uncertainty resolved at 600 dpi: Eq. (23.8.14)
% prints M^{-4}, despite the apparent dimensional inconsistency.
\subsection{Vacuum Decay}
\label{sec:23.8}

A vacuum state is stable if its scalar field expectation values are at a true minimum of the effective potential. But if the scalar field expectation values are at a local minimum that is higher than the true minimum then this vacuum will be metastable. A metastable `false' vacuum state corresponding to a local minimum will decay into the stable `true' vacuum corresponding to the true minimum by a process of barrier penetration, analogous to nuclear alpha decay or spontaneous fission. This is not a process that can be observed in our laboratories, but it has presumably occurred several times in the history of the universe as various symmetries have become spontaneously broken, so it is important to be able to calculate the rate of such false vacuum decay. As we shall now see, this calculation involves consideration of yet another extended field configuration~\hyperref[ch23-ref-6]{[6]}.

Let us concentrate on the component $\phi$ of the scalar field multiplet that acquires an expectation value $\langle\phi\rangle$ in the true vacuum. For instance, in the theory of broken chiral symmetry discussed in \hyperref[sec:19.5]{Section 19.5}, $\phi$ would be the fourth component of a chiral four-vector. It will turn out that, where barrier penetration is strongly suppressed, the other scalar fields (including those of any Goldstone bosons) do not affect the dominant suppression factor in the decay rate. For definiteness, we will take the Lagrangian density in the form
\begin{equation}
 \mathcal L=-\frac{1}{2}\partial_\mu\phi\,\partial^\mu\phi-V(\phi).
 \tag{23.8.1}\label{eq:23.8.1}
\end{equation}
We assume that the lowest-order effective potential $V(\phi)$ has a true minimum at $\phi=\langle\phi\rangle$ and a local minimum at $\phi=0$, and we adjust an additive constant in the Lagrangian density so that $V(0)=0$, in which case $V(\langle\phi\rangle)<0$. We want to calculate the rate at which the false vacuum state with scalar field expectation value zero will decay into the true vacuum state with scalar field expectation value $\langle\phi\rangle$.

The results of \hyperref[app:23.A]{Appendix A of this chapter} (Eqs.~\eqref{eq:23.A.6}, \eqref{eq:23.A.21}, and \eqref{eq:23.A.23}) show that the energy $E_0$ of the false vacuum state in which the scalar field vacuum expectation value vanishes is given by
\begin{equation}
 E_0=-\lim_{T\to\infty}\frac{1}{T}\ln\left[
 \int\exp\bigl(-S[\phi;T]\bigr)
 \prod_{\mathbf x,t}\dd\phi(\mathbf x,t)
 \right],
 \tag{23.8.2}\label{eq:23.8.2}
\end{equation}
where $S[\phi;T]$ is the Euclidean action derived from Eq.~\eqref{eq:23.8.1}
\begin{equation}
 S[\phi;T]=\int \dd^3x\int_{-T/2}^{+T/2}\dd t
 \left[
 \frac{1}{2}\left(\frac{\partial\phi}{\partial t}\right)^2
 +\frac{1}{2}(\boldsymbol\nabla\phi)^2+V(\phi)
 \right],
 \tag{23.8.3}\label{eq:23.8.3}
\end{equation}
and the integral in Eq.~\eqref{eq:23.8.2} is over all fields $\phi(\mathbf x,t)$ satisfying the conditions
\begin{equation}
 \phi(\mathbf x,T/2)=\phi(\mathbf x,-T/2)=0.
 \tag{23.8.4}\label{eq:23.8.4}
\end{equation}
The energy \eqref{eq:23.8.2} is complex; its imaginary part will give us the decay rate.

To calculate the functional integral in Eq.~\eqref{eq:23.8.2}, we look for a stationary `point' of the Euclidean action $S[\phi;T]$. The fields at which \eqref{eq:23.8.3} is stationary satisfy the field equations
\begin{equation}
 0=\frac{\delta S}{\delta\phi}
 =-\frac{\partial^2\phi}{\partial t^2}
 -\boldsymbol\nabla^2\phi+\frac{\dd V(\phi)}{\dd\phi},
 \tag{23.8.5}\label{eq:23.8.5}
\end{equation}
subject to the boundary conditions \eqref{eq:23.8.4}. Because of these boundary conditions, such a solution is known as a \emph{bounce}.

We shall look for these bounce solutions by making the ansatz that $\phi(\mathbf x,t)$ is invariant under rotations around a point $\mathbf x_0,t_0$ in four dimensions:
\begin{equation}
 \phi(\mathbf x,t)=\phi(\rho)
 \quad\text{where}\quad
 \rho\equiv\sqrt{(\mathbf x-\mathbf x_0)^2+(t-t_0)^2}.
 \tag{23.8.6}\label{eq:23.8.6}
\end{equation}
There are solutions that are not rotationally invariant in four dimensions, but these other solutions have higher values~\hyperref[ch23-ref-37]{[37]} of $S$, and so become negligible for large $T$. Using Eq.~\eqref{eq:23.8.6} in Eq.~\eqref{eq:23.8.5} yields the ordinary differential equation
\begin{equation}
 \frac{\dd^2\phi}{\dd\rho^2}+\frac{3}{\rho}\frac{\dd\phi}{\dd\rho}=V'(\phi).
 \tag{23.8.7}\label{eq:23.8.7}
\end{equation}
Strictly speaking, a solution of this form is consistent with the boundary conditions \eqref{eq:23.8.4} only when $T$ is very large compared with the characteristic time associated with $V(\phi)$, in which case we can take $T$ to be infinite in Eq.~\eqref{eq:23.8.4}, which then becomes the condition that $\phi(\rho)$ must vanish when $\rho\to\infty$. Also $\phi(\mathbf x,t)$ must be an analytic function of $\mathbf x$ near $\mathbf x=0$ at all $t$ including at $t=t_0$, so $\phi(\rho)$ is a power series in $\rho^2$ for $\rho\to0$, and in particular $\dd\phi/\dd\rho=0$ at $\rho=0$. With these conditions, Eq.~\eqref{eq:23.8.7} is the equation of motion of a particle of unit mass with `position' $\phi$ at `time' $\rho$, moving under the influence of a potential $-V(\phi)$ and a viscous force $-(3/\rho)\dd\phi/\dd\rho$, that travels from rest at some finite initial value $\phi_0$ of $\phi$ at $\rho=0$, and just reaches $\phi=0$ at $\rho\to\infty$, losing its initial `energy' $-V(\phi_0)>0$ to viscosity along the way. The Euclidean action \eqref{eq:23.8.2} for such a solution is
\begin{equation}
 B\equiv\int_0^\infty 2\pi^2\rho^3\dd\rho
 \left[
 \frac{1}{2}\left(\frac{\dd\phi}{\dd\rho}\right)^2+V(\phi)
 \right].
 \tag{23.8.8}\label{eq:23.8.8}
\end{equation}
It is crucial to determine the sign of $B$. For this purpose~\hyperref[ch23-ref-38]{[38]}, we use the same device that was used to prove Derrick's theorem in \hyperref[sec:23.1]{Section 23.1}, and consider the action \eqref{eq:23.8.8} for a modified field $\phi_R(\rho)=\phi(\rho/R)$. By rescaling the variable of integration, we find
\begin{equation}
 S[\phi_R]=\int_0^\infty 2\pi^2\rho^3\dd\rho
 \left[
 \frac{R^2}{2}\left(\frac{\dd\phi}{\dd\rho}\right)^2+R^4V(\phi)
 \right].
 \tag{23.8.9}\label{eq:23.8.9}
\end{equation}
If $\phi(\rho)$ is the solution of Eq.~\eqref{eq:23.8.7} then the action must be stationary with respect to any variation in $\phi$, so that $\dd S[\phi_R]/\dd R$ must vanish at $R=1$, and therefore
\begin{equation}
 \int_0^\infty\rho^3\dd\rho\left(\frac{\dd\phi}{\dd\rho}\right)^2
 =-4\int_0^\infty\rho^3\dd\rho\,V(\phi).
 \tag{23.8.10}\label{eq:23.8.10}
\end{equation}
It follows then that
\begin{equation}
 B=\frac{\pi^2}{2}\int_0^\infty\rho^3\dd\rho
 \left(\frac{\dd\phi}{\dd\rho}\right)^2>0.
 \tag{23.8.11}\label{eq:23.8.11}
\end{equation}
We will come back later to an explicit approximate solution for $\phi(\rho)$, but first let's consider how such solutions are to be used.

We have here not just a single one-bounce configuration at which the Euclidean action is stationary, but a continuum, characterized by the collective coordinates $\mathbf x_0$ and $t_0$. According to the results of \hyperref[sec:23.7]{Section 23.7}, we must integrate over these parameters, which, since $B$ is independent of $\mathbf x_0$ and $t_0$, yields factors of $\mathcal V$ and $T$ in a box of spatial volume $\mathcal V$. The contribution of all one-bounce configurations to the functional integral in Eq.~\eqref{eq:23.8.2} is given in one-loop order by
\begin{equation}
 \mathcal VTA\exp(-B).
 \tag{23.8.12}\label{eq:23.8.12}
\end{equation}
The coefficient $A$ is proportional\footnote{Where a continuous symmetry $G$ is spontaneously broken to a subgroup $H$, there are additional collective coordinates giving the orientation of $H$ within $G$. The integration over these parameters yields an additional factor in $A$ equal to the volume of the coset space $G/H$.} to a product $\prod_n'\lambda_n^{-1/2}$, where the $\lambda_n$ are the eigenvalues of the kernel $\delta^2S[\phi]/\delta\phi(\mathbf x,t)\delta\phi(\mathbf x',t')$, and the prime indicates that we are to omit the zero eigenvalues, which correspond to changes in the collective coordinates $\mathbf x_0$ and $t_0$. Using Eq.~\eqref{eq:23.8.10}, we can see that the second derivative of Eq.~\eqref{eq:23.8.9} with respect to $R$ is negative at $R=1$, so there is at least one (and in fact just one~\hyperref[ch23-ref-39]{[39]}) negative eigenvalue, and hence to this order $A$ is imaginary. We will not attempt to calculate $A$ here, but will just note that $A$ unlike $\exp(-B)$ does not have a dramatic dependence on the parameters of the theory, so that we can estimate it on dimensional grounds to be roughly of order $iM^{-4}$, where $M$ is some characteristic mass scale of the theory.

For large $T$ and $\mathcal V$ we can find additional stationary configurations by superimposing any number $N$ of these bounce configurations, yielding a contribution that is the $N$th power of the quantity \eqref{eq:23.8.12}, divided by $N!$ to take account of the fact that in integrating over $N$ of the $\mathbf x_0$ and $t_0$ we are summing over configurations that only differ by permutations of the $N$ identical bounces. Summing over $N$ then gives the exponential of the quantity \eqref{eq:23.8.12}, so the energy \eqref{eq:23.8.2} is just the quantity \eqref{eq:23.8.12} divided by $-T$
\begin{equation}
 E_0=-\mathcal VA\exp(-B).
 \tag{23.8.13}\label{eq:23.8.13}
\end{equation}
Since $A$ is of order $iM^{-4}$, the decay rate per volume of the false vacuum is thus of order
\begin{equation}
 \Gamma/\mathcal V\approx M^{-4}\exp(-B).
 \tag{23.8.14}\label{eq:23.8.14}
\end{equation}
Note that this is a decay rate per volume, because the decay does not occur by a change in the scalar field simultaneously everywhere in space, but by the occurrence of bubbles of true vacuum in a false vacuum background.

The result \eqref{eq:23.8.14} is chiefly useful in the case where $B$ is large, so that barrier penetration is strongly suppressed, and we can estimate the suppression factor as simply $\exp(-B)$. Fortunately, the most natural circumstance in which $B$ is large is one in which it is possible to calculate $B$ in closed form. This is the case in which the energy $V(\langle\phi\rangle)=-\epsilon$ of the true vacuum is only slightly below the zero energy of the false vacuum, but $V(\phi)$ is positive and not small between $\phi=0$ and $\phi=\langle\phi\rangle$. To minimize the Euclidean action \eqref{eq:23.8.3} in this case, we must take $\phi$ to be near $\langle\phi\rangle$ within a four-dimensional ball with a large radius $R$, at which $\phi$ drops to zero within a shell of thickness given by some length $L\simeq M^{-1}$ characteristic of the potential in the limit $\epsilon\to0$. (This is sometimes called the `thin wall approximation,' but perhaps a better term would be the big bubble approximation.) The action \eqref{eq:23.8.3} in this approximation is
\begin{equation}
 S(R)\simeq-\frac{1}{2}\pi^2R^4\epsilon+2\pi^2R^3\mathcal S,
 \tag{23.8.15}\label{eq:23.8.15}
\end{equation}
where $\mathcal S$ is a surface tension, equal to the shell contribution to the action per area. The second term on the left-hand side of Eq.~\eqref{eq:23.8.7} becomes negligible for $\rho\simeq R$, so this is now essentially a one-dimensional problem. We can therefore take the surface tension from Eq.~\eqref{eq:23.1.4}, which for a solution of the field equations is an equality, and in present notation reads
\begin{equation}
 \mathcal S=\int_0^{\langle\phi\rangle}\sqrt{2V(f)}\,\dd f.
 \tag{23.8.16}\label{eq:23.8.16}
\end{equation}
The action \eqref{eq:23.8.15} is stationary at a radius
\begin{equation}
 R\simeq3\mathcal S/\epsilon,
 \tag{23.8.17}\label{eq:23.8.17}
\end{equation}
so the action at its stationary point has the value
\begin{equation}
 B\simeq\frac{27\pi^2\mathcal S^4}{2\epsilon^3}.
 \tag{23.8.18}\label{eq:23.8.18}
\end{equation}
Note that $B$ is large for small $\epsilon$, so in this case the decay rate of the false vacuum is strongly suppressed. After the barrier has been penetrated the bubble of true vacuum will grow at the speed of light, colliding eventually with other bubbles, until all space is in the state of lowest energy.
